\chapter{Convexidad, separación y funciones sublineales}\label{cap:v3-03}
La separación geométrica necesita transformar un conjunto en una función sublineal. La funcional de Minkowski realiza ese paso, pero solo conserva la simetría cuando el conjunto la posee. Desarrollaremos las pruebas de separación antes del teorema del rango cerrado, para que este último pueda utilizar una herramienta ya demostrada. Cerraremos con una segunda construcción de extensiones por funcionales minimales y con la corrección de su fórmula auxiliar.

\section{La funcional de Minkowski}\label{lec:v3-05}
\fuentelectura{05}{Forma geométrica del teorema de Hahn--Banach}
Toda convexidad y toda desigualdad de esta sección se entienden sobre el cuerpo real. En un espacio complejo se pasa al espacio real subyacente y luego se representa una funcional real $g$ como la parte real de $\varphi(x)=g(x)-ig(ix)$. La comprobación de linealidad compleja es la misma del teorema~\ref{v3:hb-norma}.

\begin{proposicion}\label{v3:subnivel}
Para una funcional sublineal $p$ y $U_p=\{x:p(x)<1\}$, el conjunto $U_p$ es convexo y
\[
 tU_p=\{x:p(x)<t\}\quad(t>0),\qquad
 y+U_p=\{x:p(x-y)<1\}.
\]
Si $p\geq0$ y $f$ es lineal real no nula, entonces $f\leq p$ equivale a $f^{-1}(1)\cap U_p=\varnothing$. Si además $p$ es continua, $f\leq p$ implica que $f$ es continua.
\end{proposicion}
\begin{proof}
Para $x,y\in U_p$ y $0<t<1$, $p(tx+(1-t)y)\leq tp(x)+(1-t)p(y)<1$; en los extremos el resultado es inmediato. Las dos identidades se obtienen usando homogeneidad positiva y la definición de traslación.

Si $f\leq p$, un vector con $f(x)=1$ no puede satisfacer $p(x)<1$. Recíprocamente, si $f(x)>p(x)\geq0$, entonces $f(x)>0$ y $u=x/f(x)$ satisface $f(u)=1$, $p(u)<1$, contradicción. Para continuidad, dado $\eps>0$ elija $\delta>0$ tal que $\norm x<\delta$ implique simultáneamente $p(x)<\eps$ y $p(-x)<\eps$. Entonces $-p(-x)\leq f(x)\leq p(x)$ y $|f(x)|<\eps$. No se puede sustituir esta doble cota por $|f(x)|\leq p(x)$ cuando $p$ no es simétrica.
\end{proof}

\begin{teorema}[Funcional de calibración o de Minkowski]\label{v3:minkowski}
Sea $A$ convexo en un espacio normado real $X$, con $0\in\interior A$. Defina
\[
 p_A(x)=\inf\{t>0:x\in tA\}.
\]
Entonces $p_A$ es finita, no negativa, sublineal y continua. Se cumplen
\[
 \interior A=\{p_A<1\},\qquad \cl A=\{p_A\leq1\},\qquad
 \partial A=\{p_A=1\},\qquad X\setminus\cl A=\{p_A>1\}.
\]
La homogeneidad garantizada es $p_A(tx)=tp_A(x)$ para $t\geq0$; si $A=-A$, también $p_A(tx)=|t|p_A(x)$ para escalares reales.
\end{teorema}
\index{Minkowski, funcional de}\index{funcional!de calibracion@de calibración}
\begin{proof}
Elija $r>0$ con $B(0,r)\subseteq A$. Para $t>\norm x/r$, $x/t\in A$, así que $0\leq p_A(x)\leq\norm x/r$. Para $s>0$, el cambio $t=su$ en el ínfimo da $p_A(sx)=sp_A(x)$; para $s=0$ ambos miembros son cero.

Si $x\in aA$, $y\in bA$ con $a,b>0$, entonces
\[
 \frac{x+y}{a+b}=\frac a{a+b}\frac xa+\frac b{a+b}\frac yb\in A.
\]
Por ello $p_A(x+y)\leq a+b$. Elija $a,b$ en sus conjuntos definitorios a distancias arbitrariamente pequeñas de los ínfimos, para obtener subaditividad. Esta y la cota anterior dan
\[
 p_A(x)-p_A(y)\leq p_A(x-y)\leq\norm{x-y}/r,
\]
y al intercambiar $x,y$ se prueba continuidad Lipschitz.

Si $p_A(x)<1$, existe $0<t<1$ con $x=ta$, $a\in A$. Por convexidad,
\[
 x+(1-t)B(0,r)=ta+(1-t)B(0,r)\subseteq A,
\]
de modo que $x$ es interior. Si $x$ es interior y $x\ne0$, existe $s>1$ cercano a uno con $sx\in A$; luego $p_A(x)\leq1/s<1$. Para $x=0$ ya sabemos $p_A(0)=0$. Si $x\in\cl A$, la continuidad y $p_A\leq1$ en $A$ dan $p_A(x)\leq1$. Si $p_A(x)\leq1$, entonces $p_A(tx)<1$ para $0<t<1$; los puntos $tx\in\interior A$ convergen a $x$, así que $x\in\cl A$. Las fórmulas de frontera y exterior se siguen de las dos primeras. Finalmente $A=-A$ da $p_A(-x)=p_A(x)$ y con ello homogeneidad absoluta.
\end{proof}
\begin{observacion}[Una simetría que no puede suponerse]
Para $A=(-1,2)\subseteq\R$, $p_A(x)=x/2$ si $x\geq0$ y $p_A(x)=-x$ si $x<0$. Así $p_A(1)=1/2$ y $p_A(-1)=1$. La homogeneidad absoluta afirmada sin simetría en la lectura5 necesita corregirse; las demás caracterizaciones geométricas anteriores no requieren esa simetría.
\end{observacion}

\section{Separación por hiperplanos}
\begin{teorema}[Separación con interior]\label{v3:separacion-interior}
Sean $A,B$ convexos no vacíos disjuntos, con $\interior A\ne\varnothing$. Existe una funcional real continua no nula $f$ tal que
\[
 f(a)\leq f(b)\quad(a\in A,b\in B).
\]
Para $a\in\interior A$ la desigualdad es estricta.
\end{teorema}
\begin{proof}
El conjunto $C=\interior A-B$ es abierto y convexo, no vacío y no contiene cero. Elija $c_0\in C$ y ponga $U=C-c_0$, $z=-c_0$. Entonces $U$ es abierto, convexo, contiene cero y no contiene $z$. Por el teorema anterior $p_U(z)\geq1$. En $\R z$ defina $h(tz)=t$. Si $t\geq0$, $h(tz)=t\leq tp_U(z)$; si $t<0$, $h(tz)<0\leq p_U(tz)$. Hahn--Banach extiende $h$ a $f\leq p_U$. La continuidad se obtiene usando las cotas para $x$ y $-x$, y $f(z)=1$ asegura que $f\ne0$.

Para $c\in C$, $c-c_0\in U$ y $f(c)=f(c-c_0)+f(c_0)<1-1=0$. Por tanto $f(a)<f(b)$ si $a\in\interior A$, $b\in B$. Falta incluir los puntos restantes de $A$. Elija $a_0\in\interior A$ y $r>0$ con $B(a_0,r)\subseteq A$. Para $a\in A$ y $0<t\leq1$, el vector $a_t=(1-t)a+ta_0$ tiene la bola $B(a_t,tr)$ contenida en $A$; por tanto $a_t\in\interior A$. Como $a_t\to a$, la continuidad da $f(a)\leq f(b)$.
\end{proof}

\begin{teorema}[Separación estricta]\label{v3:separacion-estricta}
Si $A,B$ son convexos no vacíos disjuntos, $A$ es compacto y $B$ cerrado, existe una funcional real continua $f$ con
\[
 \sup_{a\in A}f(a)<\inf_{b\in B}f(b).
\]
En el caso complejo, las desigualdades se expresan mediante $\operatorname{Re}\varphi$ para una funcional compleja continua $\varphi$.
\end{teorema}
\begin{proof}
La función $a\mapsto\dist(a,B)$ es continua, positiva en $A$ y alcanza su mínimo $d>0$ en el compacto $A$. Elija $0<r<d$ y ponga $U=A+B(0,r)$. Es abierto, convexo y disjunto de $B$. El teorema anterior da $f\ne0$ con $f(u)\leq f(b)$ para $u\in U$, $b\in B$. Para $a\in A$, tome el supremo sobre $\norm h<r$ en $f(a+h)\leq f(b)$; resulta
\[
 f(a)+r\norm f\leq f(b).
\]
El supremo de $f(h)$ sobre esa bola es $r\norm f$ por la definición de la norma y la posibilidad de cambiar el signo de $h$. Tomando supremo en $a$ e ínfimo en $b$ se obtiene un salto de tamaño al menos $r\norm f>0$. La conversión compleja se explicó al inicio del capítulo.
\end{proof}
\index{separacion@separación!de conjuntos convexos}
\begin{corolario}[Separación absoluta de un conjunto balanceado]\label{v3:separacion-balanceada}
Sea $A\ne\varnothing$ convexo y balanceado en un espacio real o complejo: $\lambda A\subseteq A$ si $|\lambda|\leq1$. Si $x\notin\cl A$, existen $\varphi\in X^*$ y $\alpha>0$ tales que
\[
 |\varphi(a)|<\alpha<|\varphi(x)|\quad(a\in A).
\]
\end{corolario}
\begin{proof}
Separe estrictamente $\cl A$ del compacto $\{x\}$. En el caso complejo obtenga $\sup_A\operatorname{Re}\varphi<\operatorname{Re}\varphi(x)$. Como $A$ es balanceado, para cada $a$ todas sus rotaciones de módulo uno están en $A$, y $\sup_A\operatorname{Re}\varphi=\sup_A|\varphi|$. En el caso real basta cambiar el signo. Este supremo es no negativo porque $0\in A$. Elija $\alpha$ estrictamente entre el supremo y $\operatorname{Re}\varphi(x)\leq|\varphi(x)|$.
\end{proof}

\section{Ejercicios propuestos: calibración y convexidad}
\begin{ejercicio}[Lectura5, ejercicio1]\label{ej:v3-l05-1}
Calcule la funcional de calibración de la bola abierta unitaria de un espacio normado real.
\end{ejercicio}

\begin{ejercicio}[Lectura5, ejercicio2]\label{ej:v3-l05-2}
Demuestre que $f(x)=\tfrac12\norm x^2$ es convexa.
\end{ejercicio}

\begin{ejercicio}[Lectura5, ejercicio3, incisos(i)--(ii)]\label{ej:v3-l05-3}
Un espacio es estrictamente convexo si $\norm{u+v}<2$ para vectores unitarios \emph{distintos} $u,v$. (i) Pruebe la equivalencia con $\norm{(u+v)/2}<1$. (ii) Pruebe la equivalencia con la convexidad estricta de $x\mapsto\tfrac12\norm x^2$, también entendida para puntos distintos.
\end{ejercicio}


\begin{ejercicio}[Lectura5, ejercicio4: teorema del sándwich, incisos(i)--(iv)]\label{ej:v3-l05-4}
Sea $p:X\to\R$ sublineal, $S\subseteq X$ no vacío y $f:S\to\R$ con $f\leq p$ en $S$. Defina
\[
 g(x)=\inf_{s\in S,\,t\geq0}\{p(x+ts)-tf(s)\}.
\]
Pruebe: (i) $g$ es real, $g\leq p$ y el mismo ínfimo se obtiene con $t>0$; (ii) para una funcional lineal $h$, $h\leq g$ equivale a $h\leq p$ en $X$ y $h\geq f$ en $S$; (iii) si $S$ es convexo y $f$ cóncava, $g$ es sublineal; (iv) bajo esas hipótesis existe una funcional lineal entre las dos cotas. Examine también $S=\varnothing$.
\end{ejercicio}


\begin{ejercicio}[Lectura5, ejercicio5: separación de un abierto]\label{ej:v3-l05-5}
Examine si en la separación estricta puede sustituirse ``$A$ compacto'' por ``$A$ abierto'' manteniendo la misma conclusión.
\end{ejercicio}

\begin{ejercicio}[Lectura5, ejercicio6]\label{ej:v3-l05-6}
Si $H$ es un hiperplano afín cerrado real, pruebe que $H=f^{-1}(\alpha)$ para alguna $f\in X^*\setminus\{0\}$ y $\alpha\in\R$. Describa los lados y sus cierres.
\end{ejercicio}

\begin{ejercicio}[Lectura5, ejercicio7: teorema del soporte]\label{ej:v3-l05-7}
El original supone que $A$ es conexo, con interior no vacío, y pide un hiperplano cerrado por cada $x\notin\interior A$ que deje a $A$ en un semiespacio cerrado. Examine esa hipótesis y pruebe la versión válida.
\end{ejercicio}

\begin{ejercicio}[Lectura5, ejercicio8]\label{ej:v3-l05-8}
Demuestre que un convexo cerrado de un espacio normado real no nulo es la intersección de los semiespacios cerrados que lo contienen.
\end{ejercicio}


\section{Funcionales minimales y una segunda extensión}\label{lec:v3-06}
\fuentelectura{06}{Más sobre funcionales sublineales}
Las funcionales sublineales reales forman un cono ordenado por la comparación punto a punto. Una funcional $p$ es minimal si la única sublineal $q\leq p$ es $p$ misma.
\begin{proposicion}\label{v3:l06-p1}
Para una funcional sublineal $p$ son equivalentes: (i) $p$ es lineal; (ii) $p(x)+p(-x)=0$ para todo $x$; (iii) $p(x)+p(-x)\leq0$ para todo $x$.
\end{proposicion}
\begin{proof}
La linealidad implica(ii), y(ii) implica(iii). Siempre $0=p(0)\leq p(x)+p(-x)$, así que(iii) implica(ii). Si(ii) vale, entonces
\[
 p(x)\leq p(x+y)+p(-y)=p(x+y)-p(y),
\]
lo cual, junto con subaditividad, da aditividad. La homogeneidad positiva y $p(-x)=-p(x)$ dan homogeneidad para todos los escalares reales. Así(ii) implica(i), completando también la prueba propuesta en la lectura.
\end{proof}

\begin{observacion}[Corrección de la funcional auxiliar]
La fórmula original toma el ínfimo simultáneamente sobre $t\geq0$ y $w\in X$. Sin embargo,
\[
 \inf_{t\geq0,\,w\in X}\{p(x+tw)-tp(w)\}=-p(-x).
\]
La desigualdad $\geq$ procede de subaditividad y la igualdad se alcanza tomando $t=1$, $w=-x$. Para $p(x)=|x|$ en $\R$ el resultado es $-|x|$, que no es subaditivo: su valor en $1+(-1)$ es cero, mayor que $-2$. El teorema auxiliar original es falso. La corrección que conserva el método consiste en fijar la dirección $w$ antes de tomar el ínfimo.
\end{observacion}
\begin{lema}\label{v3:auxiliar-fija}
Para una dirección fija $w\in X$, defina
\[
 q_w(x)=\inf_{t\geq0}\{p(x+tw)-tp(w)\}.
\]
Entonces $q_w$ es sublineal, $-p(-x)\leq q_w(x)\leq p(x)$ y $q_w(-w)=-p(w)$.
\end{lema}
\begin{proof}
La cota inferior se obtiene de $tp(w)\leq p(x+tw)+p(-x)$, y la superior tomando $t=0$. En particular la función es real. El cambio $t=\alpha s$ demuestra homogeneidad positiva; para $x=0$ todos los términos del ínfimo son cero. Para $s,t\geq0$,
\[
 q_w(x+y)\leq p(x+y+(s+t)w)-(s+t)p(w)
 \leq[p(x+sw)-sp(w)]+[p(y+tw)-tp(w)].
\]
Los parámetros $s,t$ son independientes, así que elegir aproximantes a ambos ínfimos prueba subaditividad. Para $x=-w$, la cota inferior da $q_w(-w)\geq-p(w)$, y tomar $t=1$ da la desigualdad inversa.
\end{proof}
\begin{teorema}\label{v3:minimales}
Las funcionales sublineales minimales son exactamente las lineales. Toda funcional sublineal domina alguna funcional lineal.
\end{teorema}
\begin{proof}
Si $p$ es lineal y $q\leq p$ es sublineal, entonces $p(x)=-p(-x)\leq-q(-x)\leq q(x)\leq p(x)$, por lo que $q=p$. Si $p$ es minimal, el lema da $p=q_w$ para cada dirección $w$. Evaluar en $-w$ produce $p(-w)=-p(w)$, y la proposición anterior demuestra linealidad.

Para la existencia, considere $\mathcal P=\{q:q\text{ sublineal},q\leq p\}$, no vacío porque contiene $p$. Si $\mathcal C$ es una cadena no vacía, cada $q\in\mathcal C$ satisface $-p(-x)\leq-q(-x)\leq q(x)$; por tanto $h(x)=\inf_{q\in\mathcal C}q(x)$ es real. Es positivamente homogénea y $h(0)=0$. Dados $x,y$ y $\eps>0$, elija $q_1,q_2$ con $q_1(x)<h(x)+\eps/2$, $q_2(y)<h(y)+\eps/2$. Como la familia es una cadena, el menor de $q_1,q_2$ pertenece a ella y sirve simultáneamente para las dos cotas. Su subaditividad da $h(x+y)\leq h(x)+h(y)+\eps$. Así $h$ es sublineal y es una cota inferior en $\mathcal P$. La cadena vacía tiene la cota $p$. Zorn, aplicado al orden inverso, da un elemento minimal de $\mathcal P$; también es minimal entre todas las sublineales porque cualquier funcional menor sigue estando dominada por $p$. Es lineal por lo ya probado.
\end{proof}

\begin{proposicion}[Segunda prueba de Hahn--Banach]\label{v3:hb-minimal}
Si $f:M\to\R$ es lineal y $f\leq p$ en $M$, la función
\[
 q(x)=\inf_{m\in M}\{p(x-m)+f(m)\}
\]
es sublineal, satisface $q\leq p$ en $X$ y $q=f$ en $M$. Una funcional lineal dominada por $q$ extiende a $f$.
\end{proposicion}
\begin{proof}
La desigualdad $-f(m)=f(-m)\leq p(-m)\leq p(x-m)+p(-x)$ da $q(x)\geq-p(-x)$. Tomar $m=0$ da $q\leq p$. Si $x\in M$, entonces $f(x)=f(x-m)+f(m)\leq p(x-m)+f(m)$ y tomar $m=x$ da igualdad $q(x)=f(x)$.

La homogeneidad se obtiene cambiando $m=\alpha n$ para $\alpha>0$, y $q(0)=f(0)=0$. Para $x,y$ elija aproximantes $m,n\in M$ a los dos ínfimos; use
\[
 p(x+y-m-n)+f(m+n)\leq p(x-m)+f(m)+p(y-n)+f(n)
\]
y haga tender a cero el error de aproximación. Así $q$ es sublineal. El teorema de minimalidad produce una funcional lineal $F\leq q$. En $M$, $F\leq f$, y aplicar esa desigualdad a $-m$ da $F\geq f$. Por tanto $F|_M=f$ y $F\leq p$ en todo $X$.
\end{proof}

\section{Alcance de la separación}
La separación débil entre conjuntos no equivale a una distancia positiva entre sus imágenes. La compacidad de uno y la cerradura del otro producen esa distancia; la apertura produce desigualdades estrictas punto a punto. La funcional de Minkowski reúne estas ideas y permite verificar cada hipótesis. Su construcción, y la corrección de la auxiliar en dirección fija, muestran cómo una fórmula de ínfimos puede conservar unas propiedades y perder otras si se cambian sus parámetros.
